Definition \ref{def.hawkesprocess} leaves the kernel free.
We now fix it, and we fix it in the way that makes the process simulable and estimable
in linear time:
\begin{equation}\label{eq.exponentialKernel}
 \hawkesKernel\subscriptee(t) = \excitation\subscriptee \, e^{-\decay t},
 \qquad t \geq 0,
\end{equation}
with $\excitation\subscriptee \geq 0$ and a single decay rate $\decay > 0$ shared by every pair
of types.
We call $\excitation$ the excitation matrix and $\decay$ the decay.
From \eqref{eq.branchingMatrix},
\begin{equation}\label{eq.GammaIsAOverBeta}
 \branchingMatrix = \frac{\excitation}{\decay},
\end{equation}
so the branching matrix is dimensionless while the excitation matrix is a rate.
It is worth carrying that distinction:
$\sum_e \branchingMatrix\subscriptee$ is a count of offspring,
$\sum_e \excitation\subscriptee = \decay \sum_e \branchingMatrix\subscriptee$ is
the jump in the total intensity when one event of type $\eone$ fires,
measured per second.
The two are quoted together often enough that it is worth saying they differ by $\decay$.

\subsection{The state}

Define, for each type $e$,
\begin{equation}\label{eq.decayedCounts}
 \decayedCounts_e(t) := \sum_{j \, : \, \arrivalTimes_{j} < t} e^{-\decay (t - \arrivalTimes_{j})}.
\end{equation}
The inequality under the sum is strict.
This is not a detail: it makes $\decayedCounts$ left-continuous,
and therefore makes
\begin{equation}\label{eq.intensityFromState}
 \intensity[ ](t) = \baseIntensity + \excitation \decayedCounts(t)
\end{equation}
predictable, as Definition \ref{def.compensator} requires of an intensity.
Substituting \eqref{eq.exponentialKernel} into \eqref{eq.intensity_ordHawkes}
and separating the factor $e^{-\decay t}$ gives \eqref{eq.intensityFromState} directly.

We call $\decayedCounts$ the decayed counts.
Between events it solves $d\decayedCounts = -\decay \decayedCounts \, dt$,
and at an event of type $e$ it gains the unit vector $e_e$:
\begin{equation}\label{eq.stateDynamics}
 d\decayedCounts(t) = -\decay \decayedCounts(t) \, dt + d\multiCountingProc(t).
\end{equation}
The pair $(\multiCountingProc, \decayedCounts)$ is therefore a piecewise-deterministic Markov
process:
the whole of the past enters the future only through the $\numEventTypes$ numbers
$\decayedCounts$.
This is the property that the exponential kernel is chosen for,
and \eqref{eq.stateDynamics} is the form in which we shall use it in
Section \ref{sec.secondOrder}.

\begin{remark}[choosing the kernel is choosing the size of the state]\label{remark.stateSize}
 A single $\decay$ shared by all pairs buys a state of dimension $\numEventTypes$.
 Allowing a decay $\decay\subscriptee$ per pair would require one decayed count per pair,
 a state of dimension $\numEventTypes^2$,
 and every expression of Sections \ref{sec.stability} and \ref{sec.secondOrder} would have to be
 rewritten on it.
 What the single decay costs is that the model has one excitation timescale and cannot describe
 a market whose market orders relax quickly and whose cancellations relax slowly.
 What it buys is \eqref{eq.stateDynamics}, and with it the closed forms of this chapter.
 We take the trade throughout, and we say where it binds.
\end{remark}

\subsection{The recursion}

Let $\arrivalTimes[ ]_{n}$ and $\event[n]$ be the arrival times and labels of the ground process.
Between $\arrivalTimes[ ]_{n}$ and $\arrivalTimes[ ]_{n+1}$ the state merely decays,
so from \eqref{eq.decayedCounts},
\begin{equation}\label{eq.update_of_exponential_intensity}
 \decayedCounts(\arrivalTimes[ ]_{n+1}-)
 = \Big( \decayedCounts(\arrivalTimes[ ]_{n}-) + e_{\event[n]} \Big)
 \, e^{-\decay (\arrivalTimes[ ]_{n+1} - \arrivalTimes[ ]_{n})},
\end{equation}
and equivalently, on the intensity itself,
\begin{equation*}
 \intensity[ ](\arrivalTimes[ ]_{n+1}-)
 = \baseIntensity
 + \Big( \intensity[ ](\arrivalTimes[ ]_{n}-) + \excitation e_{\event[n]} - \baseIntensity \Big)
 \, e^{-\decay (\arrivalTimes[ ]_{n+1} - \arrivalTimes[ ]_{n})}.
\end{equation*}
Both values are needed and they are different objects.
The pre-jump value $\decayedCounts(\arrivalTimes[ ]_{n+1}-)$ is the one that enters the
intensity, by predictability;
the post-jump value $\decayedCounts(\arrivalTimes[ ]_{n+1}-) + e_{\event[n+1]}$ is the one
carried forward as the state.
Confusing them shifts every intensity by one event, which is a mistake that produces a
plausible path.

The recursion costs $O(1)$ per event.
Evaluating \eqref{eq.decayedCounts} directly at each of $n$ arrival times costs $O(n^2)$,
which is the difference between a simulation that runs and one that does not.

\subsection{The compensator in closed form}

Integrating \eqref{eq.stateDynamics} over $[0,\timeHorizon]$ with
$\decayedCounts(0) = \multiCountingProc(0) = 0$ gives
$\int_{0}^{\timeHorizon} \decayedCounts_{\eone} =
\big( \countingProc[\eone](\timeHorizon) - \decayedCounts_{\eone}(\timeHorizon) \big) / \decay$,
and therefore
\begin{equation}\label{eq.compensatorClosedForm}
 \compensator_e(\timeHorizon)
 = \baseIntensity_e \timeHorizon
 + \sum_{\eone=1}^{\numEventTypes} \frac{\excitation\subscriptee}{\decay}
 \Big( \countingProc[\eone](\timeHorizon) - \decayedCounts_{\eone}(\timeHorizon) \Big).
\end{equation}
The compensator is thus available exactly, at any time, from the state and the counts,
with no numerical integration.
Notice that $\countingProc[\eone] - \decayedCounts_{\eone}$ is continuous:
both jump by one at an arrival of type $\eone$.
So \eqref{eq.compensatorClosedForm} is insensitive to the jump convention,
and we may read it at $\timeHorizon-$ or at $\timeHorizon$ indifferently.
This is what makes the goodness-of-fit test of Section \ref{sec.simulation} cheap,
and it is the reason we compute the compensator independently of the simulator rather than
accumulating it as the path is generated:
an independent computation is a test, an accumulated one is a tautology.
